\section*{Chapter \ref{ch:hf:latency-arbitrage-and-its-defence} --- Latency Arbitrage and Its Defence}

\begin{solution}{exo:hf:latency-arbitrage-and-its-defence:1}
$2-0.5=1.5$ ticks a share, before fees.
\end{solution}

\begin{solution}{exo:hf:latency-arbitrage-and-its-defence:2}
$537/510\approx1.05$ a minute.
\end{solution}

\begin{solution}{exo:hf:latency-arbitrage-and-its-defence:3}
More than $60-45=15$ microseconds.
\end{solution}

\begin{solution}{exo:hf:latency-arbitrage-and-its-defence:4}
See \cref{prop:hf:latency-arbitrage-and-its-defence:bump}: arrival times shift by the same $d$ and keep their order. A symmetric delay, applied to the
venue's outgoing data too, still stops traders from reacting to the venue's own executions faster than the venue can route a customer's order on.
\end{solution}

\begin{solution}{exo:hf:latency-arbitrage-and-its-defence:5}
When the event falls just before the end of a batch, a taker can arrive before the boundary while the slower cancel lands after it: the taker is in
an earlier batch. The chance is about the gap between the two latencies divided by the interval, here about 2\%.
\end{solution}

\begin{solution}{exo:hf:latency-arbitrage-and-its-defence:6}
Half-way between 28.8\% and 13.4\% is 21.1\%: the guard still does it at 100 ms (18.7\%) and 300 ms (20.0\%), not at 600 ms (24.6\%).
\end{solution}

\begin{solution}{exo:hf:latency-arbitrage-and-its-defence:7}
0.314. The probability falls to 1.0\% at a 110-microsecond delay and to 0.5\% at 120: the delay must cover the provider's disadvantage plus most of
the jitter.
\end{solution}

\begin{solution}{exo:hf:latency-arbitrage-and-its-defence:8}
An order-book feed shows only the orders that succeeded; the losers of a race (orders that missed, cancels that came too late) leave no trace in it.
Races are visible only in message data with failed attempts, as Aquilina, Budish and O'Neill used, or in one's own logs of rejected cancels and
missed takes.
\end{solution}

\begin{solution}{pb:hf:latency-arbitrage-and-its-defence:1}
\begin{enumerate}
\item See \cref{def:hf:latency-arbitrage-and-its-defence:arb}; the information is public and only speed matters, unlike informed trading, which uses
information others do not have.
\item The median opportunity fell from 97 to 7 milliseconds; its median profit stayed near 0.08 index points.
\item About 537 races a day per FTSE 100 stock; modal gap 5--10 microseconds; about 22\% of volume in races; tax 0.42 basis points of all volume.
\item Message data contain failed orders and cancels, the losers; an order-book feed shows only what succeeded.
\item One provider cancel at a median 50 microseconds against three snipers at 40, lognormal jitter 0.3; a one-tick jump against a half-tick quote.
\item 0.899, and 0.449 ticks a share per race.
\item 12\%, 75\% and 99.9\%.
\item A median 7.3 microseconds, within the 5--10 microseconds reported for London.
\item See \cref{def:hf:latency-arbitrage-and-its-defence:bump}.
\item See the proof: a common shift keeps the order; with takers alone delayed the comparison becomes $x_{\text{take}}+d<x_{\text{cancel}}$.
\item 0.899 continuous and with a symmetric 350-microsecond delay; 0 with the same delay on takers only; 0.211 and 0.021 for batches of 100
microseconds and one millisecond.
\item IEX approved with a 350-microsecond delay (2016) and its venue-repriced D-Limit order (2020); EDGA's four-millisecond asymmetric delay
disapproved (2020); rescission of the trade-through rule proposed (2026).
\item The provider may reject a deal request after a hold time if the price has moved; the provider controls it, within the FX Global Code's
principles.
\item 90\% of the half-spread on a continuous book and with a symmetric delay ($0.449/0.5$), zero with an asymmetric delay, 2.1\% with a
one-millisecond batch.
\item 13.4\% and 0.402 at 1 ms, 20.0\% and 0.330 at 300 ms, 24.6\% and 0.192 at 600 ms: the guard works while it reads the leader faster than the
lead, and fails beyond it.
\item 31\% of the volume at 1 ms (5,225 shares a session against 7,550), and messages.
\item A lead guard fed as fast as it can afford, a slightly wider quote, and a preference for venues that protect quotes.
\item The delay makes their fills slower and less certain; routers that value speed of execution go elsewhere, as the SEC's record of objections
shows.
\item Venues would no longer have to route to or respect each other's protected quotes; stale quotes could be traded through instead of taken,
changing which races are worth running.
\item Liquidity providers, and through their wider spreads, everyone who trades.
\end{enumerate}
\end{solution}

\begin{solution}{iq:hf:latency-arbitrage-and-its-defence:1}
Adverse selection: the counterparty knew something you did not. Sniping: the information was public, and the counterparty was faster to act on it
than you were to cancel.
\iqlookfor{private versus public information.}
\end{solution}

\begin{solution}{iq:hf:latency-arbitrage-and-its-defence:2}
Log every cancel that arrived too late (rejects, fills between send and acknowledgement) and the fills that followed a move on another venue within
your round trip; compare the timestamps of the move, your cancel and the fill.
\iqlookfor{too-late cancels and cross-venue timestamps.}
\end{solution}

\begin{solution}{iq:hf:latency-arbitrage-and-its-defence:3}
The path from the market-data packet to the cancel on the wire: kernel bypass, feed decoding, the decision, the order encoding, the network card, the
link to the exchange; measure each stage before optimising (One Quant Book 13, chapter 1).
\iqlookfor{a measured, stage-by-stage view.}
\end{solution}

\begin{solution}{iq:hf:latency-arbitrage-and-its-defence:4}
Within a batch all orders are treated alike and priced together, so arriving earlier gains nothing; at a boundary, arriving just before rather than
just after it puts an order into an earlier auction.
\iqlookfor{discrete time and its edges.}
\end{solution}

\begin{solution}{iq:hf:latency-arbitrage-and-its-defence:5}
The firm becomes slower than its competitors on that route: its quotes there are stale for longer and its taking strategies lose races. Widen or
withdraw the affected quotes, stop the latency-sensitive strategies, and fall back to the fibre route with its latency in the risk parameters.
\iqlookfor{operational dependence on speed and a fallback.}
\end{solution}

\begin{solution}{iq:hf:latency-arbitrage-and-its-defence:6}
The minimum of $n$ exponentials with mean $s$ is exponential with rate $n/s$; it beats an independent exponential of rate $1/c$ with probability
$(n/s)/(n/s+1/c)=nc/(nc+s)$.
\iqlookfor{competing exponentials.}
\end{solution}
